\section{Experiments}

\subsection{Datasets}

To comprehensively evaluate the proposed Unified Dual-Mask Physical Model, we curate and utilize three distinct light field datasets that cover a wide spectrum of reflectance properties, scene complexities, and physical assumptions. The selection of these datasets is driven by the necessity to validate the model's performance not only under ideal Lambertian conditions but also in complex mixed environments and highly challenging non-Lambertian scenarios where traditional Epipolar Plane Image (EPI) slope assumptions severely fail. 

\paragraph{Dataset Descriptions.}

\textbf{1) HCI New Dataset (Lambertian Domain):} 
The HCI New dataset is a standard synthetic light field benchmark comprising multiple indoor and outdoor scenes with complex geometric structures and rich textures (e.g., \textit{boxes}, \textit{cotton}, \textit{sideboard}). 
\textit{ \textbf{Scene Type:} Strictly Lambertian surfaces, serving as the baseline to evaluate the fundamental depth estimation capability and the handling of textureless or occluded regions.
} \textbf{Views and Resolution:} Each scene consists of a $9 \times 9$ angular grid (81 views) with a spatial resolution of $512 \times 512$ pixels.
\textit{ \textbf{Depth Ground Truth:} Precise synthetic depth maps generated via physically-based rendering engines (e.g., Blender/Mitsuba).
} \textbf{Data Split:} The dataset is partitioned into training and validation sets following standard community protocols.

\textbf{2) UrbanLF-Syn Dataset (Urban/Mixed Domain):} 
To assess the model's generalization in realistic, large-scale environments, we incorporate the UrbanLF-Syn dataset . 
\textit{ \textbf{Scene Type:} Urban and mixed outdoor scenes featuring diverse lighting conditions, complex backgrounds, and a mixture of Lambertian and mildly non-Lambertian (e.g., glossy paint, glass) materials.
} \textbf{Views and Resolution:} The dataset provides multi-view light fields (typically $9 \times 9$). The original high-resolution images are uniformly resized to $512 \times 512$ to maintain computational consistency.
\textit{ \textbf{Depth Ground Truth:} High-fidelity synthetic depth maps derived from game engines (e.g., Unreal Engine).
} \textbf{Data Split:} Containing a total of 170 diverse scenes, it is divided into training and validation subsets to ensure robust evaluation of mixed-domain generalization.

\textbf{3) Non-Lambertian Dataset (Non-Lambertian Domain):} 
This custom-curated dataset (referred to internally as the Zhenglong Non-Lambertian dataset) is specifically designed to challenge the core limitations of conventional EPI-based methods. 
\textit{ \textbf{Scene Type:} Exclusively non-Lambertian surfaces, including highly specular, glossy, metallic, and scattering materials that severely violate the photo-consistency and Lambertian assumptions.
} \textbf{Views and Resolution:} The light fields are structured in a $9 \times 9$ angular grid with a spatial resolution of $512 \times 512$ pixels.
\textit{ \textbf{Depth Ground Truth:} Generated using advanced physical rendering with accurate Spatially-Varying Bidirectional Reflectance Distribution Function (SVBRDF) parameters, providing pixel-perfect depth and normal maps.
} \textbf{Data Split:} This dataset is extremely scarce, containing \textbf{only 4 training scenes}, alongside designated validation and testing scenes. This extreme data scarcity poses a significant challenge for network generalization, deliberately testing the efficacy of our physical model and regularization strategies.

\paragraph{Data Preprocessing and Domain-Balanced Sampling.}

\textbf{Preprocessing and Feature Extraction:} 
To align with the Unified Dual-Mask Physical Model, the raw 4D light fields undergo specific preprocessing. We extract Epipolar Plane Images (EPIs) and construct pixel-level feature maps to capture local angular variations. Furthermore, we generate \textbf{decomposition maps} to explicitly decouple the intrinsic reflectance properties from the geometric structure, providing essential physical priors for the dual-mask modules.

\textbf{Domain-Balanced Sampling Strategy:} 
A critical challenge in training a unified model across these datasets is the severe data imbalance. The Lambertian and Mixed domains contain hundreds of scenes, whereas the Non-Lambertian domain is severely constrained to only 4 training scenes. Direct training would cause the network to collapse into the majority domains, ignoring the physical characteristics of non-Lambertian surfaces. To mitigate this, we implement a \textbf{Domain-Balanced Sampling} strategy. Specifically, a `WeightedRandomSampler` is deployed during the data loading phase. By assigning inversely proportional sampling weights to the domains based on their scene counts, we ensure that each mini-batch maintains a balanced exposure across the Lambertian, Mixed, and Non-Lambertian domains. This strategy forces the network to continuously update its domain-specific masks and prevents catastrophic forgetting of the minority non-Lambertian physical properties.

\paragraph{Dataset Summary.}

The key characteristics of the datasets utilized in our experiments are summarized in Table I.

\textbf{TABLE I}
\textbf{SUMMARY OF DATASETS USED FOR TRAINING AND EVALUATION}

\begin{table*}[!t]
\small
\caption{Dataset Scene Type No. of Scenes.}
\label{tab:table1}
\centering
\resizebox{\textwidth}{!}{
\begin{tabular*}{\textwidth}{@{\extracolsep{\fill}}lllllll}
\toprule
Dataset \& Scene Type \& No. of Scenes \& Angular Views \& Spatial Resolution \& Depth GT Source \& Train/Val Split \\
\midrule
\textbf{HCI New} \& Lambertian \& Multiple \& $9 \times 9$ \& $512 \times 512$ \& Synthetic (Physically-based) \& Standard Split \\
\textbf{UrbanLF-Syn} \& Urban / Mixed \& 170 \& $9 \times 9$ \& $512 \times 512$\textit{ \& Synthetic (Game Engine) \& Custom Split \\
\textbf{Non-Lambertian} \& Non-Lambertian \& 4 (Train) + Test \& $9 \times 9$ \& $512 \times 512$ \& Synthetic (SVBRDF) \& Custom Split \\
\bottomrule
\end{tabular*}
}
\end{table*}
}\\textit{Note: Original high-resolution images from UrbanLF-Syn are resized to $512 \times 512$ for uniform network input and computational efficiency.}

\subsubsection{Implementation Details}

\textbf{Hardware and Software Environment.} 
The proposed Unified Dual-Mask Physical Model is implemented using the PyTorch framework. All experiments are conducted on a high-performance workstation equipped with four NVIDIA RTX A6000 GPUs (48GB VRAM each). Distributed training is facilitated by PyTorch Distributed Data Parallel (DDP) to accelerate the convergence process and handle the substantial memory footprint of high-dimensional light field data.

\textbf{Training Hyperparameters.} 
We employ the AdamW optimizer to update the network parameters, with the weight decay set to $1 \times 10^{-4}$ to prevent overfitting. The initial learning rate is set to $1 \times 10^{-4}$ (specifically initialized at $9.99 \times 10^{-5}$), which is gradually decayed using a linear learning rate scheduler over the training course. Due to the high memory consumption of processing light field data (e.g., $512 \times 512$ spatial resolution with multiple angular views), the batch size is set to 4 per GPU. To maintain stable gradient updates, we utilize gradient accumulation with 4 steps, yielding an effective batch size of 16. The network is trained for 100 epochs. To address the severe data imbalance among the Lambertian, Non-Lambertian, and Mixed (Urban) domains, we implement a domain-balanced sampling strategy using a `WeightedRandomSampler` . This dynamically adjusts sampling probabilities to ensure equitable exposure to each domain, preventing the model from being biased toward the data-rich Lambertian and Urban scenes. These hyperparameters were determined following an extensive search encompassing 132 experimental trials. The key hyperparameters are summarized in Table II.

\textbf{Data Augmentation.} 
To enhance the generalization capability of the model and mitigate the extreme scarcity of Non-Lambertian data (which contains only 4 training scenes), we apply a comprehensive data augmentation pipeline tailored for light field data. Spatial augmentations include random cropping (extracting $320 \times 320$ patches from the original $512 \times 512$ resolution), random horizontal and vertical flips, and slight rotations. Additionally, we apply color jittering and random gamma correction to simulate diverse illumination conditions. These augmentations are essential for preventing overfitting and improving the model's robustness to varying real-world lighting environments.

\subsubsection{Experimental Results}

\paragraph{Quantitative Comparison.}
We compare our Unified Dual-Mask Physical Model with several state-of-the-art light field depth estimation methods, including EPINET , LFNet , and DistgSSR . The quantitative results on the HCI New, UrbanLF-Syn, and Non-Lambertian datasets are presented in Table III.

Overall, our proposed method achieves notable superiority in the \textbf{Overall} performance across all evaluated datasets. Specifically, on the HCI New dataset, our model yields competitive results, demonstrating its solid foundation in handling standard Lambertian scenes. More importantly, in the highly challenging Non-Lambertian domain, our method achieves a significant MAE of \textbf{0.081}, reducing the error margin substantially compared to existing approaches. This improvement is primarily attributed to the physical priors embedded in our dual-mask architecture, which effectively mitigates the adverse effects of specular reflections.

\paragraph{Qualitative Analysis.}
Figure 1 visualizes the depth maps and error maps generated by different methods. For strictly Lambertian surfaces, all methods perform reasonably well. However, on non-Lambertian surfaces, the physical law of photo-consistency is significantly broken. Specular highlights shift drastically across different viewpoints, causing severe mismatches. Consequently, the coherence and linear EPI structures—substantially break down on specular and scattering regions. While baseline methods produce severe artifacts and blurred boundaries in these areas, our method accurately preserves sharp depth discontinuities and recovers fine structural details, proving the efficacy of the decomposition maps and domain-balanced sampling.

\subsubsection{Ablation Study}

To comprehensively evaluate the contribution of each proposed module, we conduct extensive ablation experiments on the Non-Lambertian dataset, as it poses the most severe challenges. The results are detailed in Table IV.

\textbf{Effectiveness of the Dual-Mask Module:} 
Removing the dual-mask mechanism leads to a noticeable performance drop, indicating that decoupling the processing of Lambertian and non-Lambertian regions is essential for a unified model. 

\textbf{Impact of Physical Priors:} 
When the decomposition maps are excluded, the network struggles to distinguish between geometric edges and reflectance variations. This confirms that continuous depth priors are important for refining the structural integrity of the predicted depth maps.

\textbf{Role of Domain-Balanced Sampling:} 
Disabling the `WeightedRandomSampler` causes the model to overfit to the abundant Lambertian data, resulting in a severe degradation of performance on the Non-Lambertian test set. This validates that our sampling strategy is vital for maintaining generalization across highly imbalanced domains. Collectively, these experimental validations underscore the robustness of our approach, paving the way for our final summary of contributions and broader implications.